
\subsection*{\SGARefTarget{sga7.sga7i.eviii.n4}{4}. Extensions and biextensions of smooth connected group
schemes over a field}

We recall the following lemma due to Rosenlicht [3 bis, Proposition
3].

\medskip
\noindent\underline{Lemma \SGARefTarget{sga7.sga7i.eviii.n4-1}{4.1}.}  \underline{Let $k$ be a field, and
let $X$ and $Y$ be two separable, geometrically connected $k$-schemes
of finite type, equipped respectively with $k$-rational points $e_X$
and $e_Y$.  Then every morphism}
\[
 X\times Y\longrightarrow\underline G_m
\]
\underline{which is trivial on the two subschemes $X\times e_Y$ and
$e_X\times Y$ is trivial.}

We shall use this result only when $X$ and $Y$ are group schemes over
$k$; in that case the hypotheses imply that $X$ and $Y$ are in fact
smooth over $k$.  By induction, \SGARefLink{sga7.sga7i.eviii.n4-1}{4.1} also implies the analogous result
for any number of factors.  We may therefore apply the results of VII
\SGARefLink{sga7.sga7i.evii.n1-3}{1.3} and VII \SGARefLink{sga7.sga7i.evii.n2-9}{2.9}, and obtain the following two propositions.

\medskip
\noindent\underline{Proposition \SGARefTarget{sga7.sga7i.eviii.n4-2}{4.2}.}  \underline{Let $P$ be a
smooth connected group over the field $k$.  The natural functor from
the category of group extensions of $P$ by $\underline G_m$ to the
category of line bundles on $P$ rigidified at the identity element
$e_P$ is fully faithful.  Its essential image consists of the
rigidified line bundles $\mathbb L$ on $P$ for which}
\[
 \pi^*\mathbb L\simeq
 \pr_1^*\mathbb L\otimes\pr_2^*\mathbb L,
\]
\underline{where $\pi:P\times P\longrightarrow P$ is the group law
of $P$, and $\pr_1,\pr_2$ are the two projections.}

\medskip
\noindent\underline{Proposition \SGARefTarget{sga7.sga7i.eviii.n4-3}{4.3}.}  \underline{Let $P$ and $Q$ be
two smooth connected algebraic groups over the field $k$.  The
natural functor from the category of biextensions of $(P,Q)$ by
$\underline G_m$ to the category of line bundles on $P\times Q$
birigidified with respect to the identity sections $e_P$ and $e_Q$ of
$P$ and $Q$ is fully faithful; moreover, both categories are rigid.
A birigidified line bundle $\mathbb L$ on $P\times Q$ comes from a
biextension if and only if its pullbacks to $P\times P\times Q$ and
$P\times Q\times Q$ by $\pi_P\times\id_Q$ and
$\id_P\times\pi_Q$ satisfy the ``theorem of the cube'' ([1] or
[2]).}

This last condition is simply the translation of the conditions in
VII \SGARefLink{sga7.sga7i.evii.n2-9-3}{2.9.3} b), in the explicit form given in VII \SGARefLink{sga7.sga7i.evii.n3-5}{3.5} b).

\medskip
\noindent \SGARefTarget{sga7.sga7i.eviii.n4-4}{4.4}.  I do not know whether the cube condition in \SGARefLink{sga7.sga7i.eviii.n4-3}{4.3} is
automatically satisfied.  At least it is known [2, IV 2.6] that, for
a given $\mathbb L$, there is always an integer $n\geq1$, a power of
the characteristic exponent of $k$, such that
$\mathbb L^{\otimes n}$ satisfies that condition and hence comes from
a biextension of $(P,Q)$ by $\underline G_m$.  It is also always
possible to find a purely inseparable extension $k'$ of $k$ such that
the required condition becomes true after the base extension
$k\longrightarrow k'$; after this extension, therefore, $\mathbb L$
is a birigidified line bundle arising from a biextension of $(P,Q)$
by $\underline G_m$.

In particular, when $k$ is perfect, the condition in \SGARefLink{sga7.sga7i.eviii.n4-3}{4.3} is always
satisfied.  This last fact, which formally implies the preceding
statements by a trace argument, can also be seen by applying \SGARefLink{sga7.sga7i.eviii.n4-7}{4.7}
below; cf. \SGARefLink{sga7.sga7i.eviii.n4-9}{4.9}.

